\documentclass{article}

\usepackage[T1]{fontenc}

\usepackage{booktabs,graphicx,siunitx}

\begin{document}

\section{MANCOVA}

\begin{tabular}{lllllll}
\toprule
Effect & Pillai / Wilks & approx F & df1 & df2 & p & partial $\eta$\textsuperscript{2} [95\% CI] \\
\midrule
pretest\_score & 0.17 & 23.32 & 2 & 235 & $<$.001 & 0.17 [0.10, 1.00] \\
teaching\_method & 0.48 & 37.47 & 4 & 472 & $<$.001 & 0.24 [0.18, 1.00] \\
\bottomrule
\end{tabular}

\begin{tabular}{llllll}
\toprule
DV & F & df1 & df2 & p & partial $\eta$\textsuperscript{2} [95\% CI] \\
\midrule
exam\_score & 55.07 & 2 & 236 & $<$.001 & 0.32 [0.24, 1.00] \\
retention\_score & 108.57 & 2 & 236 & $<$.001 & 0.48 [0.41, 1.00] \\
\bottomrule
\end{tabular}

\includegraphics{figures/01_mancova/figure_adjusted-means.png}

Like MANOVA, but group differences on the set of outcomes are assessed after controlling for one or more covariates. Interpret the univariate follow-ups only if the multivariate p is significant, and adjust them for the number of DVs (e.g. Bonferroni) to control familywise error. Your significance threshold ($\alpha$) is 0.05.

A MANCOVA gave a covariate-adjusted group effect, Pillai's V=.48, F(4,472)=37.47, p $<$ .001, partial $\eta$\textsuperscript{2}=.24 [.18, 1.00].

Statistical basis: Pillai's trace multivariate test (headline statistic, covariate-adjusted) (Pillai, K. C. S. (1955). "Some new test criteria in multivariate analysis." Annals of Mathematical Statistics, 26(1), 117-121.); Wilks' Lambda multivariate test (covariate-adjusted) (Wilks, S. S. (1932). "Certain generalized distributions in multivariate analysis." Biometrika, 24(3/4), 471-494.); Box's M test of homogeneity of covariance matrices (Box, G. E. P. (1949). "A general distribution theory for a class of likelihood criteria." Biometrika, 36(3-4), 317-346.); Multivariate effect size (partial $\eta$\textsuperscript{2} from the multivariate approx. F) (Ben-Shachar MS, Lüdecke D, Makowski D (2020). "effectsize: Estimation of Effect Size Indices and Standardized Parameters." Journal of Open Source Software, 5(56), 2815.); Estimated marginal (covariate-adjusted) means (Lenth RV (2024). emmeans: Estimated Marginal Means, aka Least-Squares Means. R package.).

\end{document}